% Compile with Tectonic or XeLaTeX:
% tectonic --keep-logs --keep-intermediates resume.tex
\documentclass[10pt,a4paper]{article}

\usepackage[a4paper,margin=1.45cm]{geometry}
\usepackage{fontspec}
\setmainfont[
  Path=/usr/share/fonts/truetype/lato/,
  Extension=.ttf,
  UprightFont=Lato-Regular,
  BoldFont=Lato-Bold,
  ItalicFont=Lato-Italic,
  BoldItalicFont=Lato-BoldItalic
]{Lato-Regular}
\newfontfamily\displayfont[
  Path=/usr/share/fonts/truetype/lato/,
  Extension=.ttf,
  UprightFont=Lato-Regular,
  BoldFont=Lato-Black
]{Lato-Regular}
\usepackage{polyglossia}
\setdefaultlanguage{russian}
\setotherlanguage{english}
\newfontfamily\russianfont[
  Path=/usr/share/fonts/truetype/lato/,
  Extension=.ttf,
  UprightFont=Lato-Regular,
  BoldFont=Lato-Bold
]{Lato-Regular}
\newfontfamily\cyrillicfont[
  Path=/usr/share/fonts/truetype/lato/,
  Extension=.ttf,
  UprightFont=Lato-Regular,
  BoldFont=Lato-Bold
]{Lato-Regular}
\usepackage{xcolor}
\usepackage{enumitem}
\usepackage{tabularx}
\usepackage{array}
\usepackage{hyperref}
\usepackage{titlesec}
\usepackage{ragged2e}
\usepackage{microtype}

\definecolor{ink}{HTML}{171918}
\definecolor{slate}{HTML}{4E5E63}
\definecolor{muted}{HTML}{6E7776}
\definecolor{line}{HTML}{D5D9D6}
\definecolor{accent}{HTML}{3B7180}
\definecolor{paper}{HTML}{FAFAF7}

\hypersetup{
  colorlinks=true,
  urlcolor=accent,
  pdfauthor={Rodion Islibaev},
  pdftitle={Rodion Islibaev — ML Engineer / Data Scientist},
  pdfsubject={Resume}
}
\pagestyle{empty}
\setlength{\parindent}{0pt}
\setlength{\parskip}{0pt}
\renewcommand{\arraystretch}{1.12}
\setlist[itemize]{leftmargin=1.15em,itemsep=2.5pt,topsep=3pt,parsep=0pt,partopsep=0pt}
\newcolumntype{Y}{>{\RaggedRight\arraybackslash}X}

\newcommand{\sectiontitle}[1]{%
  \vspace{10pt}%
  {\color{accent}\fontsize{8.6}{10}\selectfont\bfseries\MakeUppercase{#1}}\par
  \vspace{3pt}\color{line}\rule{\linewidth}{0.55pt}\color{ink}\vspace{5pt}
}
\newcommand{\itemtitle}[2]{%
  {\fontsize{10.5}{12}\selectfont\bfseries #1}\hfill {\color{muted}\fontsize{8.6}{10}\selectfont #2}\par
}
\newcommand{\itemmeta}[1]{%
  {\color{slate}\fontsize{8.8}{10.5}\selectfont #1}\par\vspace{2pt}
}
\newcommand{\tag}[1]{%
  \fcolorbox{line}{paper}{\strut\hspace{3pt}{\color{slate}\fontsize{7.7}{9}\selectfont #1}\hspace{3pt}}\hspace{3pt}%
}

\begin{document}
\color{ink}

{\displayfont\fontsize{25}{29}\selectfont\bfseries Родион Ислибаев}\par
\vspace{2pt}
{\color{accent}\fontsize{12}{15}\selectfont\bfseries ML Engineer \;·\; Data Scientist}\par
\vspace{7pt}
\begin{tabularx}{\linewidth}{@{}Y r@{}}
\color{slate}\fontsize{9}{11}\selectfont Уфа \;·\; открыт к стажировкам и junior-ролям \;·\; удалённо / гибрид &
\fontsize{8.8}{11}\selectfont \href{mailto:airouu.kai@gmail.com}{airouu.kai@gmail.com} \;·\; \href{https://t.me/airouuu}{@airouuu} \\
& \fontsize{8.8}{11}\selectfont \href{https://github.com/a1rou}{github.com/a1rou} \;·\; \href{https://a1rou.github.io/}{a1rou.github.io}
\end{tabularx}
\vspace{8pt}\color{line}\rule{\linewidth}{0.8pt}\color{ink}

\begin{minipage}[t]{0.655\linewidth}
\sectiontitle{Профиль}
{\fontsize{9.4}{13}\selectfont
Студент программной инженерии и участник направления Data Science в Школе 21. Собираю прикладные ML-решения: от подготовки данных и корректной валидации до работающего MVP и презентации заказчику. Основной фокус — табличные данные, классическое ML и мультимодальные AI-сервисы.
}

\sectiontitle{Избранные проекты}
\itemtitle{AI-сервис оценки ювелирных изделий}{2026 · Хакатон}
\itemmeta{Python · Streamlit · Qwen · Yandex Cloud AI Studio · Bitrix24 API}
\begin{itemize}
  \item Самостоятельно спроектировал, реализовал и развернул MVP предварительной оценки золотых украшений по фотографиям для кейса ООО «СКС Ломбард».
  \item Интегрировал мультимодальную модель Qwen для анализа типа изделия, состояния, дефектов, вставок и ориентировочного веса; нормализовал структурированные JSON-ответы.
  \item Реализовал расчёт предварительной суммы займа и выкупа по тарифной сетке, пробе, весу и состоянию; представил решение заказчику.
\end{itemize}
{\fontsize{8.4}{10}\selectfont \href{https://airouu-hackathon.streamlit.app/}{Демо} \;·\; \href{https://github.com/a1rou/hackathon}{GitHub}}\par

\vspace{7pt}
\itemtitle{Decision Trees \& Ensembles}{2026 · School 21}
\itemmeta{Python · NumPy · pandas · scikit-learn · CatBoost · LightGBM · XGBoost}
\begin{itemize}
  \item Построил временное train / validation / test-разбиение по PurchDate, исключив утечку данных из будущего.
  \item Реализовал с нуля CART, Random Forest, Gradient Boosting Decision Trees и Extra Trees; сравнил с CatBoost, LightGBM и XGBoost.
  \item Настроил модели с early stopping и оценкой по ROC-AUC и Gini; проверил переобучение на отложенной test-выборке.
\end{itemize}
{\fontsize{8.4}{10}\selectfont \href{https://github.com/a1rou/ML5_Decision_trees_School21}{GitHub}}\par

\sectiontitle{Дополнительные проекты}
{\fontsize{8.9}{11.5}\selectfont
Supervised learning, регрессия, валидация моделей, классификация и feature engineering. Полный список и ноутбуки: \href{https://github.com/a1rou}{github.com/a1rou}.
}
\end{minipage}\hfill
\begin{minipage}[t]{0.305\linewidth}
\sectiontitle{Образование}
\itemtitle{УУНиТ}{2024–2028}
\itemmeta{Бакалавриат, программная инженерия\\Уфа}
\vspace{7pt}
\itemtitle{Школа 21, Уфа}{2026–н.в.}
\itemmeta{Направление Data Science\\Проектное обучение}
\vspace{7pt}
\itemtitle{МФТИ}{Дополнительно}
\itemmeta{Алгоритмы и структуры данных\\Основы C++}

\sectiontitle{Технический стек}
\tag{Python}\tag{NumPy}\tag{pandas}\tag{SQL}\tag{MySQL}\par\vspace{4pt}
\tag{scikit-learn}\tag{CatBoost}\tag{LightGBM}\tag{XGBoost}\par\vspace{4pt}
\tag{Streamlit}\tag{Docker}\tag{Git}\tag{Linux}\tag{Bash}\tag{Jupyter}

\sectiontitle{ML-практика}
{\fontsize{8.9}{11.5}\selectfont
Feature engineering\\
Model validation\\
Hyperparameter tuning\\
Data leakage control\\
ROC-AUC · Gini · F1\\
Model explainability
}

\sectiontitle{Языки}
\begin{tabularx}{\linewidth}{@{}l Y@{}}
\fontsize{8.8}{11}\selectfont Русский & \fontsize{8.8}{11}\selectfont Родной \\
\fontsize{8.8}{11}\selectfont English & \fontsize{8.8}{11}\selectfont C1 · Advanced
\end{tabularx}
\end{minipage}

\vfill
\vspace{6pt}\color{line}\rule{\linewidth}{0.55pt}\color{muted}
\begin{center}
{\fontsize{7.7}{9}\selectfont Родион Ислибаев \;·\; Резюме \;·\; Обновлено: ноябрь 2026}
\end{center}

\end{document}
